% ECONOMICS_PROBLEM: {"id": "J42-185", "title": "Measuring employer wage-setting power", "jel_code": "J42", "source_id": "OP-0350"}
% !TeX program = xelatex
\documentclass[11pt,a4paper]{article}
\usepackage[margin=23mm]{geometry}
\usepackage{fontspec}
\setmainfont{Latin Modern Roman}
\usepackage{amsmath,amssymb,mathtools}
\usepackage{xcolor,enumitem,booktabs,longtable,array,xurl,hyperref}
\definecolor{muted}{HTML}{53636C}
\hypersetup{colorlinks=true,urlcolor=blue,linkcolor=blue}
\setlength{\parindent}{0pt}
\setlength{\parskip}{5pt}
\newcommand{\R}{\mathbb R}
\newcommand{\E}{\mathbb E}
\newcommand{\N}{\mathbb N}
\newcommand{\ind}{\mathbf 1}
\newcommand{\Var}{\operatorname{Var}}
\newcommand{\Cov}{\operatorname{Cov}}
\newcommand{\supp}{\operatorname{supp}}
\DeclareMathOperator*{\argmin}{arg\,min}
\DeclareMathOperator*{\argmax}{arg\,max}
\newcommand{\entrymeta}[3]{{\sffamily\small\color{muted}\textbf{\texttt{#1}}\quad|\quad #2\par #3\par}}
\newcommand{\Problemref}[1]{\texttt{#1}}
\newcommand{\JELref}[2]{\texttt{#2} (JEL \texttt{#1})}
\begin{document}
\section*{Measuring employer wage-setting power}
\textbf{Problem ID:} \texttt{J42-185}\quad \textbf{JEL:} \texttt{J42}\par
% BEGIN ECONOMICS STATEMENT
\label{problem:OP-0350}
{\sffamily\small\color{muted}Original subject group: Labor markets and work\par}
\label{definitions:problem:30007632}
\textbf{Audit classification:} Published research agenda. Metadata and full chapter checked. Passthrough alone is insufficient under the stated model; the counterfactual wage-component estimand is editorial.
\subsection*{Definitions and precise research target}
\textbf{Complete editorial problem.} Fix a labor market, a time interval, and a population of employer--worker matches. Let \(w_{ij}>0\) be worker \(i\)'s hourly compensation at firm \(j\), and let \(p_{ij}>0\) be the marginal revenue product of that worker's labor, holding the firm's other inputs fixed. Define the latent or model-measured wedge \(d_{ij}=\log p_{ij}-\log w_{ij}\). Let \(a_{ij}\) denote job amenities, \(s_i\) worker skills, and \(c_{ij}\) contract incentives. In a structural model \(M\), let \(w^{C}_{i}(M)>0\) be compensation when the firm's ability to restrict labor supply or bargain over rents is removed while preferences, production, and aggregate resources remain fixed. Define the competition intervention and selected counterfactual job explicitly. Compare the same baseline worker's selected counterfactual wage, even if the employer changes; report nonemployment and nonexistent positive wages separately. The market-power component is \(m_i(M)=\log w^{C}_{i}(M)-\log w_{i,j(i)}\). Given matched records, production information, and specified exogenous shocks, identify or sharply bound \(E[m_i(M)]\) over models consistent with the data. Establish which restrictions distinguish this component from amenity compensation, sorting, productivity measurement error, and incentive pay. Report sensitivity to the definition of competition and to equilibrium reallocation. The task concerns the identified set in a declared market; neither a wage wedge nor a passthrough coefficient alone constitutes the answer.

\textbf{Definitions and admissible scope.} Use a probability measure \(\mu\) on baseline workers, including a separately coded outside option \(j=0\). Production revenue \(R_j(\ell,k)\) is differentiable and \(p_{ij}=\partial R_j/\partial\ell_i\) is a latent technological marginal product, not automatically an observed variable. A structural model specifies preferences, technologies, offer sets, bargaining or wage-posting rules, resource constraints, and an equilibrium selection rule. A competition intervention \(C\) must replace those wage-setting rules with a stated price-taking or other explicit benchmark. Let \(J_i^C\) be the selected counterfactual job, and restrict the log-wage target to baseline employed workers with positive baseline and counterfactual earnings; report the mass outside this domain. For observable law \(P_O\), define \(\mathcal M(P_O)=\{M\in\mathcal M_0:P_M(O)=P_O\}\) and the identified set \(\{\int_B\log(w_i^C(M)/w_{i,j(i)}(M))\,d\mu(i\mid B):M\in\mathcal M(P_O)\}\). A bound is sharp only if every permitted value is attained or approached by such models. Fix a declared family \(\mathcal M_0\) before observing the target: its members have measurable preferences and offer sets, differentiable technologies, feasible budgets, a wage-setting game, and a selected equilibrium for both the baseline and intervention \(C\). The family includes explicit restrictions on amenities, skills, incentive contracts, productivity error, shock exogeneity, and exclusion restrictions; these restrictions are inputs to the problem. Only \(\mathcal M_0\cap\{M:P_M(O)=P_O\}\) is admissible. Choose a fixed baseline-worker subset \(B\) of positive \(\mu\)-mass on which both wages are positive and log wage changes integrable under every compared model, and use \(\mu(\cdot\mid B)\) in the log target. Define \(w_i^C\) at selected job \(J_i^C\), not at a vanished match. Nonemployment is a separate population outcome. The requested output is the identified set and its infimum/supremum, conditional on these inputs, with a proof of attainable or limiting sharpness.
\subsection*{Variants and assumption changes}
\begin{enumerate}
\item Firm-level version (editorial specialization): hold outside options fixed and impose differentiable labor supply \(L_j(w)>0\), \(L'_j(w)>0\), and profits \(R_j(L_j(w))-wL_j(w)\). At an interior optimum, \(p_j/w_j=1+1/\varepsilon_j\), where \(\varepsilon_j=w_jL'_j/L_j\). Identify \(\varepsilon_j\) using wage and employment responses rather than wage passthrough alone; this identity is already known, not a new conjecture.
\item Equilibrium version (editorial extension): allow entry, amenities, and reallocations to respond to \(C\). Replace the same-match comparison by the distribution of \((J_i^C,w_i^C)\), and jointly bound wage gains and the probability of nonemployment. This changes both the estimand and the restrictions needed to identify it.
\end{enumerate}
\subsection*{References and source audit}
\textbf{Support and access.} Metadata and full chapter checked. Passthrough alone is insufficient under the stated model; the counterfactual wage-component estimand is editorial. \textbf{Locator:} Sections 1 (pp. 658--660) and 3.1.2 (p. 684).
\begin{enumerate}\raggedright
\item\hypertarget{definitions:source:30007632:1}{} Patrick Kline. 2025. \emph{Labor Market Monopsony: Fundamentals and Frontiers}. Introduction, printed pp. 658–660; section 3.1.2, printed p. 684.
\url{https://eml.berkeley.edu/~pkline/papers/HoLE_vol6_monopsony.pdf}.
DOI: \url{https://doi.org/10.1016/bs.heslab.2025.07.007}.
\end{enumerate}

\subsection*{Common conventions: Source mathematical conventions}

These conventions apply to each entry unless explicitly replaced.

\textbf{Models and identification.} A primitive model \(m\) specifies all probability laws, feasible choices, information, equilibrium and measurement rules used by its target. The admissible class \(\mathcal M\) is fixed independently of outcomes. If \(P_O(m)\) is its observable probability law and \(\tau(m)\) the target, its identified set at observable law \(P\) is
\[
 \mathcal I_\tau(P)=\{\tau(m):m\in\mathcal M,\ P_O(m)=P\}.
\]
Point identification means that this set is a singleton. An empty set signals incompatibility of data and assumptions. Coordinate bounds need not describe the joint set. Bounds are sharp only if every retained target is attainable or arbitrarily approachable by compatible models. Finite bounds require explicit restrictions on unobserved quantities; unrestricted confounding, hidden wealth, extrapolation or arbitrary equilibrium selection can yield uninformative sets.

\textbf{Causal interventions.} A potential outcome \(Y_i(p)\) is the outcome under a complete policy regime \(p\), including timing, financing, information and market spillovers. The target population and units are fixed before treatment. With interference, use the complete assignment vector or a justified exposure mapping; individual no-interference notation alone cannot identify an economy-wide intervention. A difference of mean potential outcomes does not require the cross-world joint distribution. Conditioning on post-treatment migration, survival, enrollment or employment changes the estimand and needs separate justification.

\textbf{Statistical statements.} An estimator, confidence set, forecast or test specifies a sampling law, model class and loss/coverage criterion. Uniform \((1-\alpha)\)-coverage means \(\inf_{m\in\mathcal M}\Pr_m\{\tau(m)\in C_n\}\geq1-\alpha\), or an explicitly asymptotic version. The whole parameter space trivially has coverage, so informativeness requires a stated width, risk or power objective. A forecasting improvement is relative to a named benchmark, information set and evaluation population.

\textbf{Equilibrium and optimization.} An equilibrium comprises feasible plans, optimality, beliefs where needed, budget constraints and all market/resource clearing. The primitives or a stated benchmark define each correspondence \(\mathcal E(p;m)\). Nonemptiness is assumed or proved, never inferred just from compact policy sets. A selected-equilibrium target uses a declared selection rule; a robust target quantifies over all permitted equilibria. Use suprema or infima unless attainment is established. An \(\varepsilon\)-optimal policy is feasible and within \(\varepsilon\) of the finite optimum value. Report an empty feasible set separately.

\textbf{Welfare and accounting.} Normative utility, interpersonal normalization, social weights, numeraire, discounting and the use of public receipts are primitives. Behavioral utility need not coincide with the welfare criterion. Household welfare includes ownership income and rents once; adding profits again to compensating variation generally double-counts them. A monetary equivalent is defined by an explicit transfer and price convention, with monotonicity and range conditions. Average effects, mechanism contributions, transfers and welfare losses are different targets. Infinite sums require convergence and dynamic feasible plans require terminal or no-Ponzi conditions.

\textbf{Source versions.} A working-paper date, revision date and journal publication year are distinguished. A DOI for a journal version is not automatically the DOI of an earlier manuscript. Locators refer to the stated version and printed pages where specified. A metadata-only check does not verify a claim about a full-text conclusion. Every entry's source subsection narrows the provenance claim accordingly.
% END ECONOMICS STATEMENT
\end{document}
